% chapter4_implementation.tex
% Chapter 4: Implementation

\chapter{Implementation}

\section{Development Environment and Setup}

The system was developed on Microsoft Windows with Python~3.13 (64-bit). The full pipeline is executed by \verb|scripts/build_data_store.py| (data engineering, feature engineering, ML training, DuckDB construction), followed by \verb|scripts/process_nfhs.py| (NFHS-5 screening), before launching the dashboard via \verb|streamlit run app.py|. The test suite runs via \verb|python -m pytest -v tests|. No GPU or cloud infrastructure is required; 8\,GB RAM is recommended.

\section{Technology Stack}

Each technology was selected for a specific functional rationale. Table~\ref{tab:techstack} summarises the core dependencies.

\begin{table}[h!]
\centering
\renewcommand{\arraystretch}{1.3}
\begin{tabular}{|l|l|p{6.5cm}|}
\hline
\textbf{Technology} & \textbf{Version} & \textbf{Justification} \\
\hline
Streamlit  & $\geq$ 1.46.0 & Multi-page dashboard with native navigation and session state; eliminates separate frontend/backend. \\
DuckDB     & $\geq$ 1.0.0  & In-process columnar DB; sub-10\,ms indexed lookups on 558,211 claims; no server required. \\
scikit-learn & $\geq$ 1.4.0 & RandomForestClassifier, LogisticRegression, StandardScaler, StratifiedKFold, metrics. \\
XGBoost    & $\geq$ 2.0.0  & Gradient boosting with native \texttt{scale\_pos\_weight} for class imbalance. \\
SHAP       & $\geq$ 0.45.0 & TreeExplainer for exact Shapley-value computation on tree-based models. \\
pandas     & $\geq$ 2.2.0  & DataFrame manipulation, \texttt{groupby} aggregation, type conversion. \\
Plotly     & $\geq$ 5.20.0 & Interactive charts with hover, zoom, and download. \\
SQLite     & Built-in      & Embedded relational DB for persistent investigation case management. \\
pytest     & $\geq$ 8.0.0  & Automated data validation, query, inference, and workspace tests. \\
pyreadstat & $\geq$ 1.2.0  & Reads SPSS .sav format for NFHS-5 data ingestion. \\
\hline
\end{tabular}
\caption{Technology Stack with Versions and Justifications}
\label{tab:techstack}
\end{table}

\section{Project Folder Structure}

\begin{verbatim}
Major Project Dashboard/
+-- app.py                    # Streamlit entry point
+-- requirements.txt
+-- config/                   # settings.py (paths, thresholds, features)
+-- src/                      # data_loader, preprocessing, feature_engineering,
|                             #   predict, explainability, claims_repository,
|                             #   investigation, visualizations
+-- pages/                    # 9 Streamlit dashboard pages
+-- models/                   # rf/lr/xgb .joblib + model_metadata.json
+-- data/                     # medicare/ (raw), processed/ (Parquet),
|                             #   nfhs/, claims_fraud.duckdb, investigation.db
+-- scripts/                  # build_data_store.py, process_nfhs.py
+-- tests/                    # conftest.py + 5 test modules
\end{verbatim}

\section{Data Pipeline Implementation}

The \verb|build_data_store.py| pipeline loads four raw Medicare CSV files, concatenates inpatient and outpatient records into a unified dataset of 558,211 claims, and computes ten provider-level features via \verb|groupby().agg()| named aggregations. Five binary claim-level screening flags are applied using vectorised pandas operations. Beneficiary records are enriched with parsed dates. The pipeline concludes by writing three Parquet files (\verb|providers.parquet| $\approx$295\,KB, \verb|claims.parquet| $\approx$15\,MB, \verb|beneficiaries.parquet| $\approx$2.3\,MB) and constructing the DuckDB database ($\approx$85\,MB) with five indexed tables for sub-10\,ms lookups.

\section{Machine Learning Pipeline}

Three classifiers are trained on the ten engineered provider features. Table~\ref{tab:hyperparams} lists the hyperparameters. Five-fold stratified cross-validation is performed via \verb|StratifiedKFold(n_splits=5)| and six metrics are recorded per fold. Models are serialised with \verb|joblib.dump()| to \verb|rf_provider_model.joblib| (4.2\,MB), \verb|lr_provider_model.joblib| (2\,KB), and \verb|xgb_provider_model.joblib| (154\,KB). A holdout set (80:20 stratified split: 4,328 train, 1,082 test providers) is stored in \verb|eval_predictions.parquet| for interactive threshold evaluation.

\begin{table}[h!]
\centering
\renewcommand{\arraystretch}{1.3}
\begin{tabular}{|l|l|}
\hline
\textbf{Model} & \textbf{Hyperparameters} \\
\hline
Random Forest       & \verb|n_estimators=100, random_state=42|; default class weights \\
Logistic Regression & \verb|class_weight='balanced', max_iter=1000, random_state=42|; preceded by StandardScaler \\
XGBoost             & \verb|n_estimators=100, max_depth=4, learning_rate=0.1|, \verb|subsample=0.8|, \\
                    & \verb|colsample_bytree=0.8, scale_pos_weight=9.69, random_state=42| \\
\hline
\end{tabular}
\caption{Model Training Hyperparameters}
\label{tab:hyperparams}
\end{table}

\section{Inference and Explainability}

The \verb|ProviderRiskPredictor| class in \verb|src/predict.py| is a module-level singleton that loads all three models at import time:

\begin{verbatim}
class ProviderRiskPredictor:
    def predict_provider(self, features_dict,
        model_name="Random Forest",
        threshold=0.30, high_risk_threshold=0.70):
        X = pd.DataFrame([features_dict])[FEATURE_COLUMNS]
        prob = self._get_probability(X, model_name)
        risk_tier = ("High" if prob >= 0.70
            else "Medium" if prob >= 0.30 else "Low")
        return {"probability": float(prob),
                "risk_tier": risk_tier,
                "is_flagged": bool(prob >= threshold)}
\end{verbatim}

The \verb|explain_provider_prediction()| function in \verb|src/explainability.py| initialises a \verb|shap.TreeExplainer| on the trained Random Forest, computes Shapley values, and synthesises the top three risk-elevating and top two risk-reducing features into plain-English investigator findings.

\section{Claims Repository and Screening}

The \verb|src/claims_repository.py| module implements eleven parameterised DuckDB query functions. Dynamic \verb|WHERE| clauses are built from optional filter arguments using a column whitelist to prevent SQL injection; pagination is enforced at 500 records per page. The NFHS-5 pipeline (\verb|scripts/process_nfhs.py|) reads 724,115 records via \verb|pyreadstat|, applies six vectorised rules, and writes results to \verb|nfhs5_screened_sample.parquet| and \verb|nfhs5_rule_summary.json|.

\section{Investigation Workspace}

The \verb|src/investigation.py| module manages a SQLite database with three tables (\verb|cases|, \verb|case_notes|, \verb|audit_log|). Idempotent case creation uses \verb|INSERT ... ON CONFLICT DO NOTHING|; status updates are logged transactionally. All operations use \verb|with conn:| context managers with \verb|PRAGMA foreign_keys = ON| for ACID compliance.

\section{Frontend Implementation}

\begin{table}[h!]
\centering
\renewcommand{\arraystretch}{1.3}
\begin{tabular}{|l|p{5.5cm}|l|}
\hline
\textbf{Page} & \textbf{Key Components} & \textbf{Data Source} \\
\hline
Overview Dashboard          & KPI cards, donut/bar/trend charts, sidebar filters & DuckDB claims \\
Claim Investigation         & ILIKE search, paginated table, full claim detail    & DuckDB claims \\
Beneficiary History         & Demographics, chronological scatter timeline        & DuckDB claims \\
Provider Risk Analysis      & RF/LR/XGB probability, SHAP waterfall, findings     & Parquet + SHAP \\
Suspicious Claim Screening  & Rule-filtered table with observed values            & DuckDB CTE query \\
ML Model Performance        & CV bar chart, threshold slider, confusion matrix    & eval\_predictions.parquet \\
NFHS Medical Consistency    & Nationwide KPIs, risk distribution, sample viewer   & nfhs5\_*.parquet/json \\
Investigation Workspace     & Case CRUD, lifecycle, notes, audit trail, JSON export & SQLite workspace \\
Data Explorer               & Schema viewer, completeness \%, duplicate audit     & Parquet + DuckDB \\
\hline
\end{tabular}
\caption{Frontend Pages: Key Components and Data Sources}
\end{table}

All Plotly charts use a consistent healthcare analytics theme (Deep Teal \#0D5C75) via \verb|src/visualizations.py|. Currency amounts are converted from USD to INR at a configurable rate (default 83.0) with dual-display for traceability. All paths, thresholds, and feature column lists are centralised in \verb|config/settings.py| using \verb|pathlib.Path| for cross-platform compatibility.

\section{Automated Testing Suite}

\begin{table}[h!]
\centering
\renewcommand{\arraystretch}{1.3}
\begin{tabular}{|l|p{9cm}|}
\hline
\textbf{Test Module} & \textbf{What It Verifies} \\
\hline
\verb|test_data_validation.py|        & Parquet file integrity, DuckDB existence, ten feature columns. \\
\verb|test_claim_queries.py|          & DuckDB lookups (CLM46614), non-existent IDs, beneficiary and provider profiles. \\
\verb|test_predictions.py|            & Model loading, probability bounds ($0.0 \leq p \leq 1.0$), risk tier logic, SHAP output. \\
\verb|test_dashboard_logic.py|        & SQLite ACID safety, case idempotency, audit log, screening rules, NFHS loading. \\
\verb|test_investigation_dashboard.py| & Date filter validation and stored screening evidence consistency. \\
\hline
\end{tabular}
\caption{Automated Test Suite Coverage}
\label{tab:test-modules}
\end{table}

\section{Implementation Challenges and Solutions}

\begin{itemize}
    \item \textbf{Class imbalance (9.35\%):} Addressed via \verb|class_weight='balanced'| (LR), \verb|scale_pos_weight=9.69| (XGBoost), and ensemble voting (RF).
    \item \textbf{Provider vs.\ claim-level ground truth:} Resolved architecturally — supervised ML at provider level, transparent rule-based screening at claim level.
    \item \textbf{SHAP output format compatibility:} Handles three return formats (list of arrays, 3-D array, 2-D array) via conditional type checks.
    \item \textbf{Concurrent DuckDB access:} All query functions open DuckDB in \verb|read_only=True| mode; write access occurs exclusively during the offline build step.
\end{itemize}
